% ============================================================
% Week 7 — Feature Extraction
% Paste BEFORE \subsection{Models} in Section III.
%
% Upload to Overleaf:
%   fig7_tfidf_top_features.pdf
%   fig8_token_length_distribution.pdf
%   fig9_tsne_bert_embeddings.pdf
% ============================================================

\subsection{Feature Extraction (Week 7)}

We extract two types of features: a sparse TF-IDF representation for the
baseline model and dense BERT embeddings for the transformer models.

\subsubsection{TF-IDF Features}

Each notification is represented as a TF-IDF vector using unigrams and
bigrams, capped at 500 features with sublinear TF scaling and a minimum
document frequency of~2. The vectoriser is fitted only on the training
set to prevent data leakage, producing feature matrices of size
$519 \times 500$ (train), $109 \times 500$ (val), and $109 \times 500$ (test).

Fig.~\ref{fig:tfidf_features} shows the top~15 terms per class.
High-urgency terms relate to time pressure and security
(\textit{minutes}, \textit{code}, \textit{account}); Medium terms to
scheduling (\textit{available}, \textit{pm}); and Low terms to social
and informational content (\textit{sent}, \textit{name}).

\begin{figure}[!t]
  \centering
  \includegraphics[width=\linewidth]{fig7_tfidf_top_features.pdf}
  \caption{Top~15 TF-IDF features per urgency class by mean weight.}
  \label{fig:tfidf_features}
\end{figure}

\subsubsection{BERT Tokenisation Analysis}

We tokenise all 737 notifications with the \texttt{bert-base-uncased}
WordPiece tokeniser~\cite{devlin2019}. The mean token length is 15.4
(median 15, max 98). Table~\ref{tab:token_stats} shows that 99.7\% of
notifications fit within 64 tokens, which we adopt as
\texttt{max\_len} for fine-tuning. High and Medium messages are slightly
longer ($17.8 \pm 6.0$ and $16.8 \pm 6.2$ tokens) than Low messages
($12.7 \pm 7.4$), as shown in Fig.~\ref{fig:token_lengths}.

\begin{table}[!t]
  \centering
  \caption{BERT token-length percentiles (all 737 notifications).}
  \label{tab:token_stats}
  \begin{tabular}{lc}
    \hline
    \textbf{Statistic} & \textbf{Tokens} \\
    \hline
    Mean / Median          & 15.4 / 15 \\
    90th / 95th / 99th pct & 23 / 25 / 38 \\
    Max                    & 98 \\
    \hline
    \% $\leq$ 64 tokens    & 99.7\% \\
    \% $\leq$ 128 tokens   & 100.0\% \\
    \hline
  \end{tabular}
\end{table}

\begin{figure}[!t]
  \centering
  \includegraphics[width=\linewidth]{fig8_token_length_distribution.pdf}
  \caption{BERT token-length distribution per urgency class (histogram
    and box plot).}
  \label{fig:token_lengths}
\end{figure}

\subsubsection{Pre-trained BERT Embeddings}

To verify that BERT captures urgency-related structure before fine-tuning,
we extracted 768-d \texttt{[CLS]} vectors from \texttt{bert-base-uncased}
for all 737 notifications and projected them to 2-D with t-SNE
(perplexity~30, PCA init, 1\,000 iterations). Fig.~\ref{fig:tsne} shows
partial but visible class separation, suggesting fine-tuning has a good
starting point.

\begin{figure}[!t]
  \centering
  \includegraphics[width=0.85\linewidth]{fig9_tsne_bert_embeddings.pdf}
  \caption{t-SNE of pre-trained BERT \texttt{[CLS]} embeddings, coloured
    by urgency. Partial class separation is visible before fine-tuning.}
  \label{fig:tsne}
\end{figure}
